% ========= Themenkarte: Offene Themen (eigene Vorschläge) =========
\begin{topic}[id=open_guided,
             difficulty={Variabel},
             type={Eigenes Thema mit Exposé},
             prerequisites={Abhängig vom vorgeschlagenen Thema; erforderlich sind realistische Abgrenzung und belastbare Quellenlage},
             practice={empfohlen},
             owner={Dozententeam},
             updated={2026-04-09},
             tags={ Open Topic, Seminararbeit, Vortrag, Abstract, Abbildungen, Methodik, Evaluation, Quellen, Reproduzierbarkeit }]{Offene Themen (eigene Vorschläge)}

\TopicDescription{Eigene Themenvorschläge sind willkommen, sofern sie den Seminarvorgaben entsprechen: 
30\,min Vortrag + 10\,min Diskussion sowie eine schriftliche Ausarbeitung nach wissenschaftlichen Standards 
(mind. 4500 Worte im Hauptteil, mit Titelblatt, Abstract (de/en), Inhalts- und ggf. weiteren Verzeichnissen, 
Hauptteil inkl.\ ca.\ 25\% Abbildungen, Literaturverzeichnis, ggf.\ Anhang).}

\TopicLeitfragen{\item Ist die Forschungsfrage klar formuliert und thematisch eingegrenzt?
\item Welche Methode wird eingesetzt (Literaturvergleich, Fallstudie, Analyse, Experiment)?
\item Sind belastbare und zugängliche Primärquellen vorhanden?
\item Welcher Ergebnistyp ist geplant (z.\,B.\ Vergleich, Entwurf, Fallstudie, Evaluationsbericht)?
\item Wie wird die Mindestanforderung (4500 Worte + Abbildungen + wissenschaftliche Struktur) erfüllt?
}

\TopicScope{\item Kein unstrukturierter Ideenpitch oder bloße Meinungsäußerung.
\item Keine reine Produkt-/Tool-Vorstellung ohne Analyse oder Einordnung.
\item Keine reinen How-to-/Tutorial-Sammlungen.
\item Keine Themen ohne überprüfbare Quellenbasis oder ohne Bezug zum Seminarfach.}

\TopicPractice{Empfohlen: Kurzexposé (5–7 Sätze) mit 
Titel, Forschungsfrage, Abgrenzung, Methode, erwarteten Ergebnistypen, geplanter Abbildungsanteil 
sowie 3–5 Startquellen (peer-reviewed oder anerkannte Fachquellen). 
Dieses Exposé dient der Themenfreigabe und erleichtert die spätere Ausarbeitung.}

\TopicKeywords{Offenes Thema, Seminararbeit, Vortrag, Abstract, Abbildungen, Methodik, Evaluation, Quellen, Reproduzierbarkeit}

\nocite{seminarGuide}
\end{topic}

% ========= Exposé-Mini-Vorlage (max. 1 Seite) =========
% Kann im Katalog direkt angezeigt oder separat als Handout erzeugt werden.
% Keine Zusatzpakete erforderlich.

\newpage
\begingroup
\setlength{\parindent}{0pt}
\setlength{\parskip}{0.35em}

\begin{center}
  {\Large \textbf{Exposé-Mini-Vorlage (für eigene Themenvorschläge)}}\\
  \normalsize \emph{Bitte auf maximal 1 Seite bleiben.}
\end{center}

\textbf{Titel:} \hrulefill

\textbf{Studierende(r)/Team:} \hrulefill\quad
\textbf{Matrikel/Gruppe:} \hrulefill

\textbf{Kurzbegründung (1–2 Sätze):} Warum ist das Thema relevant (Praxis/Ingenieur-Bezug)?\\
\rule{\textwidth}{0.6pt}

\textbf{Forschungs-/Praxisfrage (genau 1 Satz):}\\
\rule{\textwidth}{0.6pt}

\textbf{Abgrenzung (Scope):}\par\smallskip
\begin{minipage}[t]{0.46\textwidth}\raggedright
\textit{In Scope (max. 3 Punkte):}
\begin{itemize}[leftmargin=*,itemsep=0.1em]
%\begin{itemize}\itemsep0.1em
  \item \;
  \item \;
  \item \;
\end{itemize}
\end{minipage}\hfill
\begin{minipage}[t]{0.46\textwidth}\raggedright
\textit{Out of Scope:}
\begin{itemize}[leftmargin=*,itemsep=0.1em]
%\begin{itemize}\itemsep0.1em
  \item \;
  \item \;
  \item \;
\end{itemize}
\end{minipage}

\textbf{Methode \& Evaluation (kurz):} \emph{z.\,B.\ Literaturvergleich, Fallstudie, Experiment/Benchmark; geplante Metriken/Baseline, Vergleichsdesign}\\
\rule{\textwidth}{0.6pt}

\textbf{Erwartete Artefakte:} \emph{z.\,B.\ Ausarbeitung, Foliensatz, Code/Notebook, Diagramme/Tabellen, Reproduzierbarkeit (README/Run-Instructions)}\\
\rule{\textwidth}{0.6pt}

\textbf{Ressourcen \& Risiken:} \emph{Daten/Hardware/Software vorhanden? Lizenzen/Zugänge? mögliche Engpässe + Plan B}\\
\rule{\textwidth}{0.6pt}

\textbf{Startquellen (3–5):}
\begin{enumerate}\itemsep0.1em
  \item 
  \item 
\end{enumerate}

\textbf{Erfüllung der Seminarvorgaben (Haken setzen):}
\begin{itemize}\itemsep0.1em
  \item 30\,min Vortrag + 10\,min Fragen sind eingeplant \hfill $\Box$
  \item Ausarbeitung mit mind.\ 4500 Worten (Hauptteil), Abstract (de/en), Verzeichnisse, \textasciitilde25\% Abbildungen \hfill $\Box$
\end{itemize}

\textit{Hinweis:} Bitte bei Verwendung von Abbildungen auf Quellen-/Bildrechte achten; KI-Assistenz (falls genutzt) transparent kennzeichnen.

\endgroup
\newpage
